\chapter{Conclusion and Future Work}
\section{Conclusion}
The project achieved its main controlled-recovery aim. A Python raw-image
workflow identifies the relevant NTFS structures, selects documented named
streams and final-cluster slack ranges, exports recovered bytes to separate
files, and records verification outcomes. The two original hidden markers
were recovered with matching expected lengths and SHA-256 values, while the
recorded working-image hash remained unchanged during recovery.

Objective 1 was achieved for the documented unencrypted laboratory image:
partition geometry and MFT location were recorded and compared with Windows
output. Objective 2 was achieved for all eight baseline-file names and lengths.
The team also confirmed matching original and recovered SHA-256 hashes
for all eight files on November 25. This later confirmation is recorded
separately from the earlier machine-generated recovery results.

Objectives 3 and 4 were achieved for the documented ADS and slack payloads.
Their original expectations and selection identities were preserved, and
later exports matched them. Objective 5 is supported by 68 passing release
tests, seven acceptance checks, and the final synthetic campaign. The
campaign recovered 50 unchanged payloads exactly and detected 50 deliberate
single-bit mismatches against the original expectations. These outcomes
demonstrate correctness on tested inputs, not immunity to every attack.

A separate 100-image detection campaign measured 100\% ADS detection with
40\% false positives and 0\% evasion, versus 60\% slack-marker detection with
0\% false positives and 40\% evasion. These fixed-rule results expose the
tradeoff between broad stream screening and a narrow plaintext signature.

Objective 6 was achieved through the tagged source package, updated README,
and small generated dataset. A fresh source copy ran without the separately
retained 8 GB images. The final report adds its own reproducible campaign
and raw results while preserving the identity of the earlier code release.
The project can therefore be considered successful within its stated
educational scope, with partial rather than universal validation.

Filename, length, original-payload hash, and source-preservation checks
support different conclusions. Each is reported separately.

\section{Future Work}
Three extensions would provide the most practical improvement:
\begin{enumerate}
\item Build a diverse Windows-created corpus with original hashes, several
allocation patterns, and independent recovery comparisons. This would
address the current concentration on one lab volume and one synthetic layout.
\item Extend strict validation consistently to the metadata CLI and implement
bounded fragmented-MFT and attribute-list handling. Each newly supported
layout should have a positive case, malformed cases, and explicit limits.
\item Extend robustness and discovery studies to independent Windows-created
positive and negative volumes. Evaluate copying, filesystem writes, and payload
removal separately from signature evasion, and compare broader detectors with
the two measured screening baselines using defined denominators.
\end{enumerate}

These extensions preserve the existing emphasis on transparent evidence.
They would broaden the tool's usefulness without requiring unsupported
claims about covert security or operational forensic readiness.

\section{Self Evaluation}
\noindent\textbf{Jeff Perez self-assigned grade: B (85/100).}

This assessment credits implementation of the parser, recovery support,
controlled slack-copy tool, verification reporting, automated tests, and
reproducible release package. Documented results demonstrate recovery of
both hidden markers and protection of the source image. Deductions reflect
limited independent real-volume validation, unsupported advanced NTFS
layouts, and the need for broader robustness and discovery experiments.
Completing the planned functionality is a substantial result, while the
remaining evaluation gaps justify a B rather than full marks.

\newpage
\noindent\textbf{Lucas Curtis self-assigned grade: A (95/100).}

This assessment credits preparation of the VirtualBox laboratory, image
exports, known-file information, the manual ADS procedure, recovery-workflow
documentation, evidence organization, and the initial final-report and
evidence-log drafts. These contributions supplied the controlled environment
and records needed to connect implementation with observed results.
Deductions acknowledge limited test-volume diversity and incomplete
historical environment details. The contribution record in Appendix C
provides the supporting development history.
